% Peekaboo-VL paper: 04_background. Included from ../peekaboo_report.tex.
\section{Background}
\label{sec:background}

A video language model does not watch a whole clip. It reads a small number
of frames, 16 or 32 in this paper, taken at even intervals. Each frame is cut into patches,
a vision encoder turns the patches into vectors, and the vectors are placed
into the language model's input next to the words of the question.
Qwen3-VL and Qwen3.5 also write a short text timestamp (for example
\texttt{<3.0 seconds>}) before each pair of frames, which is how the model
knows \emph{when} each frame was taken.

This design has two weak points on drone video. The motion between frames is
never given directly: the model must infer it from differences between frames
that may be a second apart. And with 16 frames the model sees roughly one
timestamp per second on a typical clip (two per second at 32), so it cannot separate events closer together than
that. The rest of the paper addresses both.

Our remaining design choices are standard, so we fix them here and use them
throughout. We score the four answer letters instead of generating an
answer (Sec.~\ref{sec:method}). We fine-tune with
LoRA~\cite{hu2022lora}. We measure the video's true contribution against a
blind text-only baseline (Sec.~\ref{sec:blind}), following the lesson of
Goyal et al.~\cite{goyal2017vqav2} that question sets carry answer patterns
which need no image. And we compare two systems with a paired exact McNemar
test, under success rules written down before each run.
